\chapter{The Economics of a Trading Firm}\label{ch:fm:the-economics-of-a-trading-firm}

In 2020, a year of high volatility, a listed electronic market maker's revenue rose from \$1\,517 million to \$3\,239 million, 2.13 times the year
before. Its pay rose by 2.6\% and its communication and data bill by 2.1\%. The same cost base that had left it with a pre-tax loss of \$116 million
in 2019 left it with a pre-tax profit of \$1\,383 million in 2020. A listed quantitative asset manager, paid on the assets it manages and on the
performance it delivers, lived a different cycle over the same years: its revenue rose when its trend-following funds did well and fell when they
did not, whatever the volatility index did, and it paid a large share of each marginal dollar out in bonuses. This chapter reads both firms, and a
bank's markets division, through one standard set of lines, and asks of each the same two questions: which costs are fixed, and how far can
revenue fall before profit is gone.

\section{The income statement of a trading firm}

Every trading firm, whatever it calls its lines, spends its revenue in the same four places: on costs that grow with the business it does (exchange
and clearing fees, payments for order flow, the financing of the positions it holds), on pay, on the fixed machinery of the firm (technology,
market data, premises, administration), and on what is left for its owners and the tax authority.

\begin{definition}[Net trading revenue, volume-driven cost]\label{def:fm:the-economics-of-a-trading-firm:ntr}
A trading firm's \emph{volume-driven costs}\index{volume-driven cost} are the costs that grow in proportion to the business it does: exchange,
clearing and brokerage fees, payments for order flow, and the financing of its positions and of the margin behind them. Its \emph{net trading
revenue}\index{net trading revenue} is its trading and commission revenue, including interest and dividends earned on positions, less its
volume-driven costs.
\end{definition}

\omterm{def:fm:the-economics-of-a-trading-firm:ntr}{Net trading revenue} is the market maker's revenue capture (One Quant Book 11, chapter 1)
summed over a year's volume, and the right top line for comparing firms: two firms that earn the same after fees are the same business, whatever
their gross revenue. A market maker that holds large inventories financed by securities lending reports both the interest it earns and the
interest it pays; netting them is what makes its revenue comparable with an asset manager's fees.

\begin{definition}[Fixed cost base]\label{def:fm:the-economics-of-a-trading-firm:fixed}
A firm's \emph{fixed cost base}\index{fixed cost base} is the part of its costs that does not move with its revenue within a year: fixed pay,
technology and communication, market data, premises, administration and depreciation. Its \emph{operating margin}\index{operating margin} is
operating profit, \omterm{def:fm:the-economics-of-a-trading-firm:ntr}{net trading revenue} less pay less the other fixed costs, as a share of \omterm{def:fm:the-economics-of-a-trading-firm:ntr}{net trading revenue}.
\end{definition}

\begin{omfigure}
\begin{tikzpicture}[font=\footnotesize,>=stealth,
  box/.style={draw,rounded corners,minimum width=3.2cm,minimum height=0.62cm,align=center},
  cost/.style={box,fill=omThm!10},total/.style={box,fill=omDef!12}]
\node[total] (rev) at (0,4.2) {revenue};
\node[cost] (vol) at (5.2,4.2) {volume-driven costs\\{\scriptsize fees, clearing, order flow, financing}};
\node[total] (ntr) at (0,2.8) {net trading revenue};
\node[cost] (pay) at (5.2,2.8) {pay\\{\scriptsize fixed part + variable part}};
\node[cost] (fix) at (5.2,1.4) {other fixed costs\\{\scriptsize technology, data, premises, admin}};
\node[total] (op) at (0,1.4) {operating profit};
\node[box] (below) at (0,0) {one-offs, debt interest, tax};
\node[total] (ni) at (5.2,0) {net income};
\draw[->] (rev) -- (ntr); \draw[->] (ntr) -- (op); \draw[->] (op) -- (below); \draw[->] (below) -- (ni);
\draw[->,omThm] (vol.west) -- ++(-1.1,0) |- (ntr.east);
\draw[->,omThm] (pay.west) -- ++(-1.1,0) |- (op.north east);
\draw[->,omThm] (fix.west) -- (op.east);
\node[font=\scriptsize,omThm,anchor=west] at (7.7,4.2) {scale with volume};
\node[font=\scriptsize,omThm,anchor=west] at (7.7,2.8) {partly scale with revenue};
\node[font=\scriptsize,omThm,anchor=west] at (7.7,1.4) {fixed within a year};
\end{tikzpicture}

\omcaption{The standard lines of this book. Every filer's reported lines are mapped onto them (\texttt{firm.firmecon}); the costs on the right are
taken off in order, and what each one does when revenue moves is written beside it.}\label{fig:fm:the-economics-of-a-trading-firm:lines}
\end{omfigure}

The mapping is a table from each filer's reported lines to the standard ones (\cref{lst:fm:the-economics-of-a-trading-firm:map}). Three filers of
three business models run through the chapter: an electronic market maker's annual reports on Form 10-K for 2019--2025, a listed quantitative asset
manager's annual results for 2020--2025, and a bank's markets segment for 2023--2025. Their figures, in \$ millions as filed, are committed as small
derived tables with the chapter's code.

\omcode{code/desk/01-the-economics-of-a-trading-firm/python/fm_econ.py}{23}{40}{The mapping tables: each standard line is a signed sum of a
filer's reported lines.}\label{lst:fm:the-economics-of-a-trading-firm:map}

\begin{example}[The market maker's 2025 net trading revenue]\label{ex:fm:the-economics-of-a-trading-firm:virtu}
Virtu Financial's 2025 income statement reports trading income of \$2\,436.7 million, commissions and technology services of \$617.0 million,
interest and dividends income of \$508.8 million, and among its expenses brokerage, exchange, clearance fees and payments for order flow of
\$769.8 million and interest and dividends expense of \$647.4 million. Its \omterm{def:fm:the-economics-of-a-trading-firm:ntr}{net trading revenue} is $2\,436.7+617.0+508.8-769.8-647.4=\$2\,145.3$
million; the \omterm{def:fm:the-economics-of-a-trading-firm:ntr}{volume-driven costs} take $1\,417.2/3\,562.5=39.8\%$ of the revenue before other income. Pay (\$528.1 million) and the other fixed
costs (communication and data \$249.2 million, operations and administration \$97.9 million, depreciation and amortisation \$64.4 million, \$411.5
million in all) leave an operating profit of \$1\,205.7 million, a margin of 56.2\% on \omterm{def:fm:the-economics-of-a-trading-firm:ntr}{net trading revenue}.
\end{example}

The firm itself reports the same \$2\,145.3 million under another name, adjusted net trading income, and calls it a \omterm{def:fm:the-economics-of-a-trading-firm:nongaap}{non-GAAP measure}. Most listed
trading firms publish such measures, because the accounting standards were not written for their businesses; the reader's task is to find the
reconciliation and check that the adjustments go both ways.

\begin{definition}[Non-GAAP measure]\label{def:fm:the-economics-of-a-trading-firm:nongaap}
A \emph{non-GAAP measure}\index{non-GAAP measure} is a figure of performance or position that a company publishes alongside its statutory
accounts but that excludes amounts the accounting standards include, or includes amounts they exclude: adjusted revenue, core profit, adjusted
earnings. Securities regulators require it to be labelled, defined and reconciled to the nearest statutory figure.
\end{definition}

\begin{dated}{2026-09}{The rules on non-GAAP and alternative measures}\label{dat:fm:the-economics-of-a-trading-firm:rules}
In the United States, Regulation G (17 CFR 244.100--102) requires a registrant that discloses a non-GAAP financial measure to present the most
directly comparable GAAP measure with it and a quantitative reconciliation of the two; Item 10(e) of Regulation S-K applies the same to filings with
the SEC. In the European Union, ESMA's Guidelines on Alternative Performance Measures (ESMA/2015/1415, 5 October 2015) require such measures to be
defined, given meaningful labels and reconciled to the most directly reconcilable line of the financial statements. Man Group, for example, states
that its ``core'' measures are alternative performance measures and reconciles them to statutory profit.
\end{dated}

\section{Costs that scale and costs that do not}

The three firms spend their revenue in very different proportions (\cref{fig:fm:the-economics-of-a-trading-firm:shares}). The market maker gives
four dollars in ten to exchanges, clearing houses, brokers and lenders before it pays anyone; the asset manager's \omterm{def:fm:the-economics-of-a-trading-firm:ntr}{volume-driven costs}, its asset
servicing, are five in a hundred, but it pays almost half its revenue to its staff; the bank segment reports no volume-driven line of its own,
because its fees sit inside other operating expenses, and those are larger than its pay.

\begin{omfigure}
\begin{tikzpicture}
\begin{axis}[width=11cm,height=5.4cm,xbar stacked,bar width=13pt,xmin=0,xmax=100,xlabel={\small share of revenue (\%)},
  ytick={0,1,2},yticklabels={market maker,asset manager,bank markets},y dir=reverse,ymin=-0.6,ymax=2.6,
  tick label style={font=\footnotesize},table/col sep=comma,grid=major,grid style={gray!20},
  legend style={font=\scriptsize,at={(0.5,-0.32)},anchor=north,/tikz/every even column/.append style={column sep=8pt}},legend columns=4]
\addplot[fill=omThm!55,draw=omThm] table[y=k,x=volume]{figdata/desk/01-the-economics-of-a-trading-firm/shares.csv};
\addplot[fill=omProp!55,draw=omProp] table[y=k,x=comp]{figdata/desk/01-the-economics-of-a-trading-firm/shares.csv};
\addplot[fill=gray!35,draw=gray] table[y=k,x=fixed]{figdata/desk/01-the-economics-of-a-trading-firm/shares.csv};
\addplot[fill=omDef!55,draw=omDef] table[y=k,x=profit]{figdata/desk/01-the-economics-of-a-trading-firm/shares.csv};
\legend{volume-driven costs,pay,other fixed costs,operating profit}
\end{axis}
\end{tikzpicture}

\omcaption{Where each firm's 2025 revenue went, in per cent of revenue before other income: an electronic market maker (Virtu Financial), a listed
quantitative asset manager (Man Group, core measures) and a bank's markets segment (Goldman Sachs, Global Banking \& Markets, net of credit
provisions). Data: Forms 10-K for 2025 and Man Group's results for 2025, through \texttt{fm\_econ.latest\_shares}.}\label{fig:fm:the-economics-of-a-trading-firm:shares}
\end{omfigure}

The split between costs that scale and costs that do not decides what a bad year does to a firm. Write \omterm{def:fm:the-economics-of-a-trading-firm:ntr}{net trading revenue} as $\mathrm{Rev}$ (net of
\omterm{def:fm:the-economics-of-a-trading-firm:ntr}{volume-driven costs}), the fixed part of the costs, fixed pay included, as $C^{\mathrm{fix}}$, and let variable pay take a share $b$ of each
dollar of net revenue. Operating profit is
\[
\Pi(\mathrm{Rev})=(1-b)\,\mathrm{Rev}-C^{\mathrm{fix}} .
\]

\begin{proposition}[Break-even fall and operating leverage]\label{prop:fm:the-economics-of-a-trading-firm:fall}
If $\Pi(\mathrm{Rev})>0$, operating profit reaches zero when \omterm{def:fm:the-economics-of-a-trading-firm:ntr}{net trading revenue} falls by the fraction
\[
x^\ast=1-\frac{C^{\mathrm{fix}}}{(1-b)\,\mathrm{Rev}},
\]
and the elasticity of operating profit to \omterm{def:fm:the-economics-of-a-trading-firm:ntr}{net trading revenue}, the operating leverage, is
\[
\frac{d\ln\Pi}{d\ln\mathrm{Rev}}=\frac{(1-b)\,\mathrm{Rev}}{(1-b)\,\mathrm{Rev}-C^{\mathrm{fix}}}=\frac1{x^\ast}.
\]
For given total costs at the current revenue, a larger variable share $b$ raises the break-even fall and lowers the operating leverage.
\end{proposition}

\begin{proof}
Solve $(1-b)(1-x)\mathrm{Rev}=C^{\mathrm{fix}}$ for $x$; differentiate $\ln\Pi$. For the comparison, hold total costs
$K=C^{\mathrm{fix}}+b\,\mathrm{Rev}$ fixed: then $\Pi=\mathrm{Rev}-K$ does not depend on $b$, while $(1-b)\mathrm{Rev}$, the numerator of the
leverage, falls as $b$ rises, and $x^\ast=\Pi/((1-b)\mathrm{Rev})$ rises.
\end{proof}

Operating leverage is Book 11's proposition for a market maker's volume (chapter 1 there) written for the whole firm: there the fixed costs faced
volume, here they face net revenue after every \omterm{def:fm:the-economics-of-a-trading-firm:ntr}{volume-driven cost}. The one new ingredient is pay. A firm that pays a fixed share of every marginal
dollar in bonuses shares its bad years with its staff; a firm whose pay is mostly salary does not.

\begin{example}[A firm of 100]\label{ex:fm:the-economics-of-a-trading-firm:small}
A firm has \omterm{def:fm:the-economics-of-a-trading-firm:ntr}{net trading revenue} of 100, fixed pay of 15, variable pay of 30\% of net revenue and other fixed costs of 35: operating profit is
$100-15-30-35=20$. With pay that flexes, $C^{\mathrm{fix}}=50$ and $b=0.3$: the break-even fall is $1-50/70=28.6\%$ and the operating leverage 3.5.
Were the same 45 of pay all salary, $C^{\mathrm{fix}}=15+30+35=80$ and $b=0$: the break-even fall is $1-80/100=20\%$ and the
operating leverage 5. Flexible pay turns a firm that
loses money after a 20\% fall into one that survives a 28\% fall.
\end{example}

\begin{method}[Reading a firm's cost structure from its filings]\label{met:fm:the-economics-of-a-trading-firm:read}
\begin{enumerate}
\item Map every reported line onto the standard lines and check that they add up to reported pre-tax profit, the one-off items apart.
\item Compute \omterm{def:fm:the-economics-of-a-trading-firm:ntr}{net trading revenue}, pay and other fixed costs for every year available (five at least).
\item Regress pay on \omterm{def:fm:the-economics-of-a-trading-firm:ntr}{net trading revenue} across the years: the slope estimates $b$, the intercept fixed pay. Read the slope's standard error before
the slope.
\item Take the other fixed costs at their latest level, and compute $x^\ast$ and the operating leverage at the latest year's revenue, with the
estimated $b$ and with $b=0$.
\end{enumerate}
\end{method}

\section{Pay: the largest line, and how much of it is variable}

\begin{definition}[Compensation ratio]\label{def:fm:the-economics-of-a-trading-firm:comp}
A firm's \emph{compensation ratio}\index{compensation ratio} is its total pay, salaries, bonuses, deferred awards and payroll taxes, as a share of
its net revenue over the same period. Firms that publish it state their own definition of net revenue.
\end{definition}

A \omterm{def:fm:the-economics-of-a-trading-firm:comp}{compensation ratio} is not a cost structure. The asset manager's ratio was 40\% in 2021 and 2022 and 50\% in 2023 (Man Group's own figures,
pay over core net revenue before asset-servicing costs); its stated policy is a ratio ``between 40\% and 50\% of core net revenue, depending on the
mix and level of revenue''. The ratio rose in 2023 because revenue fell by more than pay did: a ratio moves against revenue whenever part of pay is
fixed. The slope of pay against revenue across years says how much of pay is in fact variable
(\cref{fig:fm:the-economics-of-a-trading-firm:pay}).

\begin{omfigure}
\begin{tikzpicture}
\begin{axis}[width=10.5cm,height=5.8cm,xlabel={\small net trading revenue (\$ million)},ylabel={\small pay (\$ million)},
  xmin=800,xmax=2400,ymin=330,ymax=720,tick label style={font=\footnotesize},scaled ticks=false,
  xtick={800,1200,1600,2000,2400},xticklabels={800,1200,1600,2000,2400},
  grid=major,grid style={gray!20},table/col sep=comma,legend style={font=\scriptsize,at={(0.5,-0.3)},anchor=north,/tikz/every even column/.append style={column sep=8pt}},legend columns=2]
\addplot[only marks,mark=*,omThm] table[x=nr,y=comp]{figdata/desk/01-the-economics-of-a-trading-firm/pay_virtu.csv};
\addplot[omThm,thick] table[x=nr,y=fit]{figdata/desk/01-the-economics-of-a-trading-firm/pay_virtu.csv};
\addplot[only marks,mark=square*,omDef] table[x=nr,y=comp]{figdata/desk/01-the-economics-of-a-trading-firm/pay_man.csv};
\addplot[omDef,thick,dashed] table[x=nr,y=fit]{figdata/desk/01-the-economics-of-a-trading-firm/pay_man.csv};
\legend{market maker (2019--2025),its fit: slope 0.05,asset manager (2020--2025),its fit: slope 0.29}
\end{axis}
\end{tikzpicture}

\omcaption{Pay against \omterm{def:fm:the-economics-of-a-trading-firm:ntr}{net trading revenue}, one point per year, with the least-squares line of each firm. The market maker's pay hardly moves
with revenue (slope 0.048, standard error 0.045); the asset manager's rises by 29 cents per dollar (standard error 0.094). Data: Virtu Financial
Forms 10-K 2019--2025, Man Group results 2020--2025, through \texttt{fm\_econ.pay\_fit}.}\label{fig:fm:the-economics-of-a-trading-firm:pay}
\end{omfigure}

The market maker's pay was between \$376 million and \$394 million in every year from 2019 to 2023, while its \omterm{def:fm:the-economics-of-a-trading-firm:ntr}{net trading revenue} ranged from
\$975 million to \$2\,271 million; the slope is 0.048 with a standard error of 0.045, indistinguishable from zero. Its 2025 pay, \$528 million, sits
far above the line, the first sign that it had begun to change. The asset manager's slope, 0.29 with a standard error of 0.094, says that about
three dimes of each marginal dollar of its revenue went to pay. The bank segment's three years give a slope of 0.20; three points make a line but
not an estimate. The asset manager also reports the split directly for 2021: fixed compensation \$208 million, variable \$388 million.

\begin{remark}[What six points can say]\label{rem:fm:the-economics-of-a-trading-firm:six}
A slope from six or seven annual points has four or five degrees of freedom and moves a great deal when one year is dropped: without 2020, the
asset manager's slope falls from 0.29 to 0.13, with a standard error of 0.13 (exercise 7). Growth in headcount, an acquisition or a change in the
mix of fees shifts pay for reasons that have nothing to do with the year's revenue. The regression is a first reading of the cost structure, to be
checked against whatever the firm says about its fixed and variable pay, never a measurement of it.
\end{remark}

\begin{center}
\small
\begin{tabular}{@{}lrrr@{}}
\toprule
2025 & market maker & asset manager & bank markets\\
\midrule
\omterm{def:fm:the-economics-of-a-trading-firm:ntr}{net trading revenue} (\$ million) & 2\,145.3 & 1\,325 & 41\,075\\
pay / \omterm{def:fm:the-economics-of-a-trading-firm:ntr}{net trading revenue} & 24.6\% & 50.9\% & 26.8\%\\
\omterm{def:fm:the-economics-of-a-trading-firm:fixed}{operating margin} & 56.2\% & 32.3\% & 42.8\%\\
\omterm{def:fm:the-economics-of-a-trading-firm:ntr}{net trading revenue} per head (\$ million) & 2.09 & 0.77 & --\\
pay per head (\$ million) & 0.51 & 0.39 & --\\
\bottomrule
\end{tabular}
\end{center}

\section{Capital, financing and return on equity}

\begin{definition}[Return on equity]\label{def:fm:the-economics-of-a-trading-firm:roe}
A firm's (or a business's) \emph{return on equity}\index{return on equity}, ROE, is its net income over a period divided by its average
shareholders' equity over the period; for a division of a larger firm, the equity is the share of the group's equity the firm attributes to it.
\end{definition}

The market maker earned net income, non-controlling interests included, of \$912.3 million in 2025 on average total equity of \$1\,730 million:
an ROE of 52.7\%. In 2023, the quiet year, the same calculation gives 17.3\%. The bank segment reports a return on average common equity of 16.4\%
for 2025, on attributed equity of \$79.7 billion. The two numbers do not measure the same thing. A market maker's equity is what its owners have
left in the firm; it is bound by the capital rules for brokers and investment firms (chapter 2), and the balance sheet it finances with its
brokers' and lenders' money is far larger than the equity. A bank's segment equity is an allocation, set by whichever of the group's capital
constraints binds (chapter 5), and the group holds more of it than its markets business could ever lose in a year, by rule.

\begin{dated}{2026-09}{Three firms in their 2025 filings}\label{dat:fm:the-economics-of-a-trading-firm:three}
\textbf{Virtu Financial} (Form 10-K for 2025): total revenue \$3\,632.1 million; volume-driven lines \$1\,417.2 million (brokerage, exchange,
clearance fees and payments for order flow \$769.8 million, interest and dividends expense \$647.4 million); employee compensation \$528.1 million;
adjusted net trading income, a \omterm{def:fm:the-economics-of-a-trading-firm:nongaap}{non-GAAP measure}, \$2\,145.3 million; about 1\,027 employees in February 2026.
\textbf{Man Group} (results for 2025): core net revenue \$1\,398 million, of which net management fees \$1\,077 million and performance fees \$281
million; core compensation \$675 million, a \omterm{def:fm:the-economics-of-a-trading-firm:comp}{compensation ratio} of 48\% (47\% in 2024); core profit before tax \$407 million; 1\,719 staff at the end
of the year. \textbf{Goldman Sachs, Global Banking \& Markets} (Form 10-K for 2025): net revenues \$41\,453 million; compensation and benefits
\$11\,025 million; other operating expenses \$12\,476 million; pre-tax earnings \$17\,574 million; average common equity \$79\,748 million; return on
average common equity 16.4\% (13.8\% in 2024, 11.3\% in 2023).
\end{dated}

\begin{remark}[Financing is a volume-driven cost]\label{rem:fm:the-economics-of-a-trading-firm:financing}
Between 2021 and 2025 the market maker's interest and dividends expense rose from \$140 million to \$647 million, and its interest and dividends
income from \$75 million to \$509 million: higher rates raised both sides of the financing of its inventory. Net, the change was small; gross, it
more than doubled the reported \omterm{def:fm:the-economics-of-a-trading-firm:ntr}{volume-driven costs}. That is why this book nets financing into \omterm{def:fm:the-economics-of-a-trading-firm:ntr}{net trading revenue}, and why a firm's reported
revenue can rise with interest rates while its business does not grow.
\end{remark}

\section{The cycle: volatility, volume and revenue}

A market maker's net revenue is volume times capture (One Quant Book 11), and both rise with volatility. Regressing the market maker's \omterm{def:fm:the-economics-of-a-trading-firm:ntr}{net trading
revenue} on the year's mean of the VIX index gives a slope of \$52 million per index point (standard error 33) and a correlation of 0.58 over seven
years. The asset manager's net revenue has a correlation of $-0.27$ with the same index: its revenue follows its assets and its funds'
performance, not market volatility. Over the seven years the market maker's pay and fixed costs moved within a narrow band while revenue swung
by a factor of 2.3 (\cref{fig:fm:the-economics-of-a-trading-firm:virtu}).

\begin{omfigure}
\begin{tikzpicture}
\begin{axis}[width=11cm,height=5.4cm,xlabel={\small year},ylabel={\small \$ million},xmin=2018.6,xmax=2025.4,ymin=0,ymax=2500,
  xtick={2019,2020,2021,2022,2023,2024,2025},xticklabel style={/pgf/number format/1000 sep={}},scaled ticks=false,
  tick label style={font=\footnotesize},grid=major,grid style={gray!20},table/col sep=comma,
  legend style={font=\scriptsize,at={(0.5,-0.3)},anchor=north,/tikz/every even column/.append style={column sep=8pt}},legend columns=4]
\addplot[omDef,thick,mark=*] table[x=year,y=nr]{figdata/desk/01-the-economics-of-a-trading-firm/virtu_years.csv};
\addplot[omProp,thick,mark=square*] table[x=year,y=comp]{figdata/desk/01-the-economics-of-a-trading-firm/virtu_years.csv};
\addplot[gray,thick,mark=triangle*] table[x=year,y=fixed]{figdata/desk/01-the-economics-of-a-trading-firm/virtu_years.csv};
\addplot[omThm,thick,dashed,mark=diamond*] table[x=year,y=profit]{figdata/desk/01-the-economics-of-a-trading-firm/virtu_years.csv};
\legend{net trading revenue,pay,other fixed costs,operating profit}
\end{axis}
\end{tikzpicture}

\omcaption{The market maker, 2019--2025: \omterm{def:fm:the-economics-of-a-trading-firm:ntr}{net trading revenue} (revenue less fees, order-flow payments and financing), pay, other fixed costs and
operating profit (before intangible amortisation, debt interest, one-off items and tax). Data: Virtu Financial Forms 10-K for 2019--2025, through
\texttt{fm\_econ.table}.}\label{fig:fm:the-economics-of-a-trading-firm:virtu}
\end{omfigure}

Applying \cref{met:fm:the-economics-of-a-trading-firm:read} at each firm's latest year gives its break-even fall and operating leverage, with pay
flexing at the estimated slope and with pay held at its latest level:

\begin{center}
\small
\begin{tabular}{@{}lrrrr@{}}
\toprule
 & \multicolumn{2}{c}{pay flexes} & \multicolumn{2}{c}{pay held fixed}\\
 & break-even fall & leverage & break-even fall & leverage\\
\midrule
market maker (2025) & 59.0\% & 1.69 & 56.2\% & 1.78\\
asset manager (2025) & 45.6\% & 2.19 & 32.3\% & 3.10\\
bank markets (2025) & 53.3\% & 1.88 & 42.8\% & 2.34\\
\bottomrule
\end{tabular}
\end{center}

\begin{omfigure}
\begin{tikzpicture}
\begin{axis}[width=10.5cm,height=5.6cm,xlabel={\small change in net trading revenue (\%)},ylabel={\small operating profit (latest year $=1$)},
  xmin=-60,xmax=40,ymin=-0.4,ymax=2,tick label style={font=\footnotesize},grid=major,grid style={gray!20},table/col sep=comma,
  legend style={font=\scriptsize,at={(0.5,-0.3)},anchor=north,/tikz/every even column/.append style={column sep=8pt}},legend columns=3]
\addplot[omThm,thick] table[x=change,y=virtu]{figdata/desk/01-the-economics-of-a-trading-firm/cycle.csv};
\addplot[omDef,thick,dashed] table[x=change,y=man]{figdata/desk/01-the-economics-of-a-trading-firm/cycle.csv};
\addplot[omProp,thick,dotted] table[x=change,y=gsgbm]{figdata/desk/01-the-economics-of-a-trading-firm/cycle.csv};
\draw[gray] (axis cs:-60,0) -- (axis cs:40,0);
\legend{market maker,asset manager,bank markets}
\end{axis}
\end{tikzpicture}

\omcaption{Operating profit after a change in \omterm{def:fm:the-economics-of-a-trading-firm:ntr}{net trading revenue}, relative to the latest year's, with pay flexing at each firm's estimated slope.
Each line crosses zero at its break-even fall: 59\% for the market maker, 46\% for the asset manager, 53\% for the bank segment. Data:
\texttt{fm\_econ.profit\_curves} on the 2025 filings.}\label{fig:fm:the-economics-of-a-trading-firm:cycle}
\end{omfigure}

The market maker, the firm with the most rigid pay, is the most robust of the three at its 2025 revenue, because its margin is widest: a high margin
is a large distance to break-even, however it was earned. The asset manager's flexible pay is what keeps it alive in a bad year: held fixed, its
pay would put break-even at a 32\% fall instead of 46\%. In 2019 the market maker's margin was 21.8\%, and its break-even fall at that year's
revenue, with the same pay slope, was 22.9\%, against 59.0\% in 2025. A firm's robustness is a property of its cost base \emph{and} of where it stands in the cycle.

\section{Tutorial: three firms through one set of lines}\label{tut:fm:the-economics-of-a-trading-firm:tut}

\begin{tutorial}
\textbf{Goal.} Read three firms' filings through one chart of accounts and measure each one's cost structure. \textbf{End state:} the table of
break-even falls and \cref{fig:fm:the-economics-of-a-trading-firm:cycle}.
\begin{tutsteps}
\item \textbf{The data.} \texttt{data/desk/filings\_virtu.csv}, \texttt{filings\_man.csv} and \texttt{filings\_gsgbm.csv} hold each filer's lines,
in \$ millions, with the filing and page recorded in \texttt{data/desk/LICENSES.md}; \texttt{vix\_annual.csv} holds the index's yearly mean,
written by \texttt{fm\_fetch\_vix.py} (the only step that needs the network).
\item \textbf{Map.} \texttt{fm\_econ.table(firm)} maps each year onto the standard lines with \texttt{firm.firmecon.statements} and the tables of
\cref{lst:fm:the-economics-of-a-trading-firm:map}. Check that the market maker's \omterm{def:fm:the-economics-of-a-trading-firm:ntr}{net trading revenue} reproduces its own non-GAAP figure,
\$2\,145.3 million in 2025 and \$1\,597.7 million in 2024.
\item \textbf{Fit pay.} \texttt{fm\_econ.pay\_fit(firm)} regresses pay on \omterm{def:fm:the-economics-of-a-trading-firm:ntr}{net trading revenue} (\cref{fig:fm:the-economics-of-a-trading-firm:pay}).
\item \textbf{Break-even.} \texttt{fm\_econ.structure(firm, flex)} builds the cost structure and computes $x^\ast$ and the leverage
(\cref{lst:fm:the-economics-of-a-trading-firm:cs}); \texttt{profit\_curves()} traces \cref{fig:fm:the-economics-of-a-trading-firm:cycle}.
\omcode{code/firm/firmecon/firm_firmecon.py}{100}{126}{The cost structure from a firm's own years, and the break-even fall and operating
leverage of \cref{prop:fm:the-economics-of-a-trading-firm:fall}.}\label{lst:fm:the-economics-of-a-trading-firm:cs}
\item \textbf{The cycle.} \texttt{fm\_econ.vix\_fit(firm)} regresses \omterm{def:fm:the-economics-of-a-trading-firm:ntr}{net trading revenue} on the index's yearly mean.
\end{tutsteps}
\end{tutorial}

\textbf{What to change next.} Drop one year at a time and watch each slope move (exercise 7); add a fourth filer from its annual report, writing
only a new mapping table.

\section{Build: the firm's economics as code}\label{bld:fm:the-economics-of-a-trading-firm:firmecon}

\begin{build}
\bfield{Purpose} One chart of accounts for every business model in the book, so that later chapters (the partnership, the platform, the fund, the
bank desk, the entry plan) state their economics in the same lines.
\bfield{Interface} \texttt{firm.firmecon}: \texttt{Statement(year, net\_revenue, comp, other\_fixed, volume\_cost, headcount, equity)} with
\texttt{operating\_profit}; \texttt{map\_lines(row, mapping)}, \texttt{statements(rows, mapping)}; \texttt{ratios(st)}; \texttt{fit\_line(x, y)}
with standard errors; \texttt{CostStructure(fixed, var\_share)}, \texttt{cost\_structure(sts, flex\_pay)}; \texttt{profit},
\texttt{breakeven\_fall}, \texttt{operating\_leverage}, \texttt{cycle}.
\bfield{Rules} A mapping names every reported column it uses, with its sign; nothing is estimated from fewer than three years; operating leverage is
undefined at or below break-even and raises an error there.
\bfield{Acceptance tests} \texttt{code/firm/firmecon/tests/}: a mapping reproduces a hand-built statement; the fitted line is exact on exact data;
profit is zero after the break-even fall; the leverage equals the numerical elasticity; flexible pay lowers leverage and raises the break-even fall
at equal total cost.
\bfield{Stretch} A parser from a filer's XBRL facts to the mapping's columns; a cycle scenario in which \omterm{def:fm:the-economics-of-a-trading-firm:ntr}{volume-driven costs}, net revenue and
variable pay all move with a volatility path.
\end{build}

\begin{omsources}
\item Virtu Financial, Inc., Forms 10-K for 2019 to 2025 (SEC EDGAR), and the SEC's XBRL company facts for the firm.
\item Man Group plc, Results for the years ended 31 December 2021, 2023 and 2025.
\item The Goldman Sachs Group, Inc., Form 10-K for 2025, segment results.
\item Cboe Global Markets, VIX index daily history (derived annual means only).
\item US Code of Federal Regulations, 17 CFR Part 244 (Regulation G); ESMA, Guidelines on Alternative Performance Measures, ESMA/2015/1415.
\end{omsources}

\section{Exercises}

\begin{exercise}[$\star$]\label{exo:fm:the-economics-of-a-trading-firm:1}
From the market maker's 2025 lines in \cref{ex:fm:the-economics-of-a-trading-firm:virtu}, compute its \omterm{def:fm:the-economics-of-a-trading-firm:ntr}{net trading revenue} and the share of its
revenue before other income taken by \omterm{def:fm:the-economics-of-a-trading-firm:ntr}{volume-driven costs}.
\end{exercise}

\begin{exercise}[$\star$]\label{exo:fm:the-economics-of-a-trading-firm:2}
The asset manager paid \$675 million on core net revenue of \$1\,398 million in 2025 and \$678 million on \$1\,696 million in 2022. Compute its
\omterm{def:fm:the-economics-of-a-trading-firm:comp}{compensation ratio} in each year, as it defines it, and say why the ratio rose although pay fell.
\end{exercise}

\begin{exercise}[$\star$]\label{exo:fm:the-economics-of-a-trading-firm:3}
The bank segment's net earnings to common were \$13\,117 million in 2025 on average common equity of \$79\,748 million. Check its reported \omterm{def:fm:the-economics-of-a-trading-firm:roe}{return on
equity}.
\end{exercise}

\begin{exercise}[$\star\star$]\label{exo:fm:the-economics-of-a-trading-firm:4}
For the firm of \cref{ex:fm:the-economics-of-a-trading-firm:small}, compute operating profit after a 25\% fall in \omterm{def:fm:the-economics-of-a-trading-firm:ntr}{net trading revenue} with pay
flexing, and with all pay fixed.
\end{exercise}

\begin{exercise}[$\star\star$]\label{exo:fm:the-economics-of-a-trading-firm:5}
From 2019 to 2020 the market maker's \omterm{def:fm:the-economics-of-a-trading-firm:ntr}{net trading revenue} rose from \$974.7 million to \$2\,271.4 million and its operating profit from \$212.4
million to \$1\,502.8 million. Compute both percentage changes and their ratio, and compare the ratio with its operating leverage in 2019.
\end{exercise}

\begin{exercise}[$\star\star$]\label{exo:fm:the-economics-of-a-trading-firm:6}
In \cref{fig:fm:the-economics-of-a-trading-firm:pay}, which of the market maker's points lies furthest above its line, and what two readings of it
would you check against the filing?
\end{exercise}

\begin{exercise}[$\star\star\star$]\label{exo:fm:the-economics-of-a-trading-firm:7}
\emph{Coding.} Refit the asset manager's pay on its \omterm{def:fm:the-economics-of-a-trading-firm:ntr}{net trading revenue} without 2020. What are the slope and its standard error, and what does the
change say about the estimate?
\end{exercise}

\begin{exercise}[$\star\star\star$]\label{exo:fm:the-economics-of-a-trading-firm:8}
\emph{Find the flaw.} ``The market maker earned 52.7\% on its equity in 2025 and the bank's markets business 16.4\%: market making is more than
three times as good a business as banking.''
\end{exercise}

\section{Problem: Three Firms, One Bad Year}

\begin{problem}[{Weekend problem --- three firms, one bad year}]\label{pb:fm:the-economics-of-a-trading-firm:1}
A board asks how each of three listed businesses would survive a year in which its \omterm{def:fm:the-economics-of-a-trading-firm:ntr}{net trading revenue} falls by a third. You have their filings.

\medskip\noindent\textbf{Part I --- The lines.}
\begin{enumerate}
\item Define \omterm{def:fm:the-economics-of-a-trading-firm:ntr}{net trading revenue} and \omterm{def:fm:the-economics-of-a-trading-firm:ntr}{volume-driven costs}, and say why financing belongs among the latter.
\item Compute the market maker's \omterm{def:fm:the-economics-of-a-trading-firm:ntr}{net trading revenue} in 2025, and check it against the figure it publishes.
\item Give each firm's 2025 operating profit and \omterm{def:fm:the-economics-of-a-trading-firm:fixed}{operating margin} on \omterm{def:fm:the-economics-of-a-trading-firm:ntr}{net trading revenue}.
\item Give the market maker's and the asset manager's \omterm{def:fm:the-economics-of-a-trading-firm:ntr}{net trading revenue} and pay per head.
\item State what Regulation G and ESMA's guidelines require of a non-GAAP or alternative measure.
\end{enumerate}

\medskip\noindent\textbf{Part II --- The cost structure.}
\begin{enumerate}[resume]
\item Write operating profit as a function of \omterm{def:fm:the-economics-of-a-trading-firm:ntr}{net trading revenue} with fixed costs $C^{\mathrm{fix}}$ and a variable share $b$.
\item Derive the break-even fall and the operating leverage, and show that one is the reciprocal of the other.
\item Give each firm's estimated pay slope $b$ with its standard error, and say which estimates you trust.
\item Give each firm's $C^{\mathrm{fix}}$ with pay flexing.
\item Why does the asset manager's \omterm{def:fm:the-economics-of-a-trading-firm:comp}{compensation ratio} move against its revenue?
\end{enumerate}

\medskip\noindent\textbf{Part III --- The bad year.}
\begin{enumerate}[resume]
\item Give each firm's break-even fall with pay flexing and with pay held fixed.
\item Give each firm's operating leverage in both cases.
\item What share of its 2025 operating profit does each firm keep after a 30\% fall, with pay flexing?
\item Which firm survives a fall of a third, and which only because its pay flexes?
\item What was the market maker's \omterm{def:fm:the-economics-of-a-trading-firm:fixed}{operating margin} in 2019, and its break-even fall at that year's revenue?
\end{enumerate}

\medskip\noindent\textbf{Part IV --- The verdict.}
\begin{enumerate}[resume]
\item Give the slope and correlation of the market maker's and the asset manager's \omterm{def:fm:the-economics-of-a-trading-firm:ntr}{net trading revenue} against the VIX index's yearly mean.
\item Compute the market maker's \omterm{def:fm:the-economics-of-a-trading-firm:roe}{return on equity} in 2025 and in 2023.
\item Why do its \omterm{def:fm:the-economics-of-a-trading-firm:roe}{return on equity} and the bank segment's measure different things?
\item State the \emph{named result}: each firm's break-even fall and operating leverage with pay flexing and held fixed.
\item In two sentences, what should the board of a firm with rigid pay and a thin margin do first?
\end{enumerate}
\end{problem}

\section{Interview questions}

\begin{interviewq}[$\star$ \iqroles{trader, researcher}]\label{iq:fm:the-economics-of-a-trading-firm:1}
What is a \omterm{def:fm:the-economics-of-a-trading-firm:comp}{compensation ratio}, and why can it rise in a year when pay falls?
\end{interviewq}

\begin{interviewq}[$\star$ \iqroles{trader}]\label{iq:fm:the-economics-of-a-trading-firm:2}
Why is a market maker's revenue higher in volatile years? Name the two channels.
\end{interviewq}

\begin{interviewq}[$\star\star$ \iqroles{risk}]\label{iq:fm:the-economics-of-a-trading-firm:3}
Two firms have the same revenue and the same total costs this year; one pays 30\% of each marginal dollar as bonus, the other only salaries. Which
has the higher operating leverage, and which survives a larger fall?
\end{interviewq}

\begin{interviewq}[$\star\star$ \iqroles{trader, risk}]\label{iq:fm:the-economics-of-a-trading-firm:4}
A firm has net revenue of 100, fixed costs of 50 and variable costs of 20\% of net revenue. What is its operating profit after a 40\% fall in
revenue?
\end{interviewq}

\begin{interviewq}[$\star\star$ \iqroles{bank}]\label{iq:fm:the-economics-of-a-trading-firm:5}
A proprietary market maker reports a \omterm{def:fm:the-economics-of-a-trading-firm:roe}{return on equity} above 50\%; a bank's markets division reports 16\%. Give three reasons they cannot be
compared directly.
\end{interviewq}

\begin{interviewq}[$\star\star\star$ \iqroles{researcher}]\label{iq:fm:the-economics-of-a-trading-firm:6}
You regress a firm's pay on its revenue over six years and find a slope of $0.29\pm0.09$. How much do you trust it, and what else could produce it?
\end{interviewq}
